% Arquivo de configuracao do documento
% Define os dados do autor e do trabalho

% O template aceita dois campos de autor (A e B). Para os cinco integrantes,
% o campo B recebe os demais nomes separados por \autsep, que vira virgula
% na pagina do artigo e quebra de linha na capa (ver o fim deste arquivo).
\newcommand{\autsep}{, }

\autorNomeA{Erick A.}
\autorSobrenomeA{Bonifácio}

\autorNomeB{Augusto Juiz Ribeiro\autsep Pedro Henrique Mendes da Silva\autsep Wesley Batista Luz\autsep João Paulo da Costa}
\autorSobrenomeB{Santos}

% Contato do(s) autor(es): informar o email
\emailAutorA{d2023005736@unifei.edu.br}
\emailAutorB{d2023004523@unifei.edu.br, d2025010682@unifei.edu.br, d2023008237@unifei.edu.br, d2025009583@unifei.edu.br}

% O titulo do trabalho
\titulo{Perfis de evasão no ensino superior a partir de grafos de similaridade entre estudantes}

% Professor da disciplina (campo "orientador" do template)
\orientadorNome{Rafael de Magalhães Dias}
\orientadorSobrenome{Frinhani (\textit{frinhani@unifei.edu.br})}

% Sem coorientador
\coorientadorNome{}
\coorientadorSobrenome{}

\txtOrientador{Professor}
\txtCoorientador{Coorientador}

% Ano de publicacao do trabalho
\ano{2026}

% Mes da entrega da versao final
\mes{Novembro}

% Local: nao altere
\local{Itajubá}

% Palavras chave (3 a 5)
\palavraChaveA{Grafos}
\palavraChaveB{Evasão estudantil}
\palavraChaveC{Detecção de comunidades}
\palavraChaveD{Redes neurais em grafos\ignorespaces}
\palavraChaveE{Dados mistos}

% Banca e data de defesa: nao se aplicam a projeto de disciplina
% (a folha de aprovacao nao e gerada; ver o fim deste arquivo)
\membroBancaA{}
\membroBancaB{}
\membroBancaC{}
\membroBancaD{}
\membroBancaE{}
\dataDefesa{}
\pdfAprovacao{}

% Pacotes adicionais
\usepackage{booktabs}
\usepackage{tikz} % diagrama do fluxo (Figura 1)
\usetikzlibrary{positioning}
% Paleta unica das figuras (mesmas cores de src/cores.py)
\definecolor{corLeiden}{HTML}{534AB7}
\definecolor{corSage}{HTML}{E08A00}
\definecolor{corEvadiu}{HTML}{C0392B}
\definecolor{corFormou}{HTML}{27AE60}

% Mais espaco para figuras de largura total no topo das paginas
\setcounter{dbltopnumber}{3}
\renewcommand{\dbltopfraction}{0.85}
\renewcommand{\dblfloatpagefraction}{0.6}

% ---------------------------------------------------------------------------
% Ajustes para relatorio de disciplina (sem alterar imc.cls):
% - o limite do projeto e de 9 paginas incluindo capa e contracapa, entao
%   ficha catalografica, folha de aprovacao e pagina em branco (proprias de TCC)
%   nao sao geradas;
% - o tipo do trabalho na capa passa a ser o projeto da disciplina;
% - o cabecalho das paginas impares lista os sobrenomes dos cinco autores.
% Para voltar ao comportamento original do template, apague este bloco.
% ---------------------------------------------------------------------------
\makeatletter
\renewcommand{\tipo}{Projeto da disciplina SMAC03 -- Grafos}
\renewcommand{\maketitle}{%
    \hypersetup{%
        pdftitle={\@titulo},
        pdfauthor={Erick A. Bonifacio, Augusto Juiz Ribeiro, Pedro Henrique Mendes da Silva, Wesley Batista Luz, Joao Paulo da Costa Santos},
        pdfsubject={SMAC03 - Grafos - UNIFEI},
    }%
    {\renewcommand{\autsep}{\\}\twocolumn[\criaCapaDura]}%
    \clearpage
    \criaPaginaInicialArtigo%
}
\fancyhead[LO]{\slshape Bonifácio, Ribeiro, Silva, Luz e Santos}
\makeatother
